\scanpage{6}
\unsure{the|thh} resolution; this \deleted{\unsure{t}}should not be hard whenever we have existence of resolutions in the strong sense of Hironaka, for instance in char.$0$ of if $\dim A=2$, \inserted{\handwritten{so that two resolutions can be dominated by \unsure{-|a} third one.}} I should point out that, in any case, the algebraic structure thus defined on $\operatorname{Pic}(U)$ depends in an essential way on the choice of the lifting of the residue field; an instructive example on which to study the dependance of this structure on the lifting is the one where $A$ is the completion of the local ring at the origin of the projecting cone of an elliptic curve (here $Q$ is isomorphic to the el liptic curve in question, for the \uline{given} lifting of $k$ ...).

I have to correct an inaccuracy that has slipped in, when defining $Q$ in the \unsure{wase|case} when $k$ is not separably closed, so as to get (10). One will then have to replace \unsure{thh|the} constant group scheme $D_k$ by the twisted constant group scheme $\underline{D}$, deduced by descent from the constant group-scheme occuring after passage from $k$ to the separable closure. \deleted{In terms of \deleted{\unsure{former ring|the ring} $B_o=H^{\unsure{o|0}}(X_o,\underline{O}_{X_o})$}\inserted{the components $X_o^i$ of $X_o$,} this can be defined as \unsure{thh|the} direct image, under the morphism $\operatorname{Spec}(B_{o\,\mathrm{sep}})\quad\operatorname{Spec}\unsure{(k)|kk)}$, of the constant sheaf $\underline{Z}$ on \deleted{\unsure{Sp}} the spectrum of the largest separable sub-algebra $B_{o\,\mathrm{sep}}$ of $B_o$} This can be made explicit, in terms of the separable algebraic closures \deleted{of} $k_i$ of $k$ in the function \unsure{fields|fielde} of the irreducible components $X_o^i$ of $X_o$, as the \deleted{product of the schemes} direct image of the constant sheaf $\underline{Z}$ under the projection $\operatorname{Spec}(\handwritten{\prod}k_i)\handwritten{\longrightarrow}\operatorname{Spec}(k)$

To study the inverse system $(P_n)$, use the fact (SGA 6 XII) that the \deleted{\unsure{s}}transition morphisms are affine, which gives a sense to $P=\underset{\handwritten{\longleftarrow}}{\lim}P_n$ as an actual group-scheme. It's Lie algebra is the inverse limit of the Lie algebras Hf the $P_n$, which are $H^1(X_n,\underline{O}_{X_n})$, therefore we get
% archive sheet number: 5
